\section{La madurez de la infraestructura de inferencia: separación PD, EP grande y colaboración nube-borde}

\subsection{Por qué la separación PD se popularizó tan rápido}
La primera palabra clave que el invitado mencionó repetidamente fue la separación PD, es decir, Prefill-Decode Disaggregation. La razón por la que en 2025 pasó rápidamente de ser un concepto de investigación a una arquitectura de producción es que las dos etapas del workload de inferencia tienen necesidades de hardware completamente distintas: Prefill consume más cómputo, mientras que Decode consume más ancho de banda y gestión del KV cache. Atar ambas etapas a un mismo conjunto de máquinas equivale, en esencia, a desperdiciar recursos.

\begin{importantbox}{El valor de ingeniería de la separación PD}
La separación PD permitió, por primera vez, que los clústeres de inferencia contaran con una forma de orquestación de recursos con mayor granularidad:
\begin{itemize}
  \item los nodos de Prefill pueden optimizarse para un alto poder de cómputo;
  \item los nodos de Decode pueden optimizarse para un ancho de banda alto y una gran capacidad de caché;
  \item el sistema en su conjunto puede escalar de forma independiente con mayor facilidad según los cambios de tráfico;
  \item la velocidad de transferencia de tecnología entre la comunidad de código abierto y los despliegues industriales se aceleró notablemente.
\end{itemize}
\end{importantbox}

Zhang Mingxing (张明星) mencionó que, cuando apenas comenzó el proyecto relacionado el año pasado, todavía había pocos usuarios; para 2025, impulsados por la colaboración de la comunidad, casi todos los proveedores principales habían adoptado arquitecturas similares. Esto ilustra bien una tendencia: cuando el workload se vuelve suficientemente grande, la optimización a nivel de arquitectura llega a un consenso en la industria más rápido que la optimización local de kernels.

\subsection{EP grande e inferencia con MoE: otra línea de expansión de capacidad}
Otra línea principal de la infraestructura de inferencia es el Expert Parallelism a gran escala. A medida que los modelos MoE siguen expandiéndose, el costo del enrutamiento de expertos y de la comunicación entre tarjetas es cada vez mayor; la Infra ya no solo tiene que responder «si el modelo cabe o no», sino también «cómo dividir a los expertos», «cómo evitar que la comunicación entre nodos arrastre la latencia» y «cómo equilibrar el servicio entre el rendimiento y la latencia de cola».

\begin{warningbox}{La dificultad de la Infra de MoE no es simplemente añadir más tarjetas}
Cuantos más Experts hay, más complejos se vuelven el enrutamiento, el balanceo de carga y la transferencia de activaciones entre dispositivos. Si la estrategia de programación y la topología de despliegue no están bien diseñadas, un EP grande puede, por el contrario, hacer que el sistema salga perdiendo en comunicación y fragmentación.
\end{warningbox}

\subsection{El lado de la nube y el lado del borde no son el mismo tipo de problema de optimización}
Otro punto práctico del panel es distinguir con claridad entre la inferencia en la nube y la inferencia en el borde. Los dos proyectos de código abierto en los que trabaja el invitado se orientan, uno hacia la optimización de la inferencia en clústeres distribuidos del lado de la nube, y el otro hacia la optimización de la inferencia heterogénea CPU-GPU del lado del borde. Esto corresponde a dos objetivos de sistema completamente distintos: el lado de la nube busca el rendimiento total, la agrupación de recursos y la eficiencia multiinquilino; el lado del borde se preocupa más por la latencia de una sola máquina, el consumo de energía y la privacidad local.

\begin{table}[H]
\centering
\begin{tabular}{p{2.8cm}p{4.0cm}p{4.5cm}}
\toprule
\textbf{Forma de despliegue} & \textbf{Objetivo principal} & \textbf{Puntos típicos de atención de la Infra} \\
\midrule
Clúster en la nube & Rendimiento, escalado, utilización multiinquilino & Separación PD, EP grande, comunicación entre nodos, orquestación de servicios \\
Borde & Baja latencia, bajo consumo de energía, privacidad, disponibilidad local & Heterogeneidad CPU-GPU, poda de operadores, uso de memoria, adaptación de dispositivos \\
Modo híbrido & Equilibrio entre costo y experiencia & Localización de la ruta caliente, migración de la ruta fría a la nube, sincronización de estado \\
\bottomrule
\end{tabular}
\caption{Aunque ambos se llaman «optimización de inferencia», los objetivos de optimización del lado de la nube y del lado del borde no son los mismos}
\end{table}

\subsection{Resumen de la sección}
El avance más importante de la Infra de inferencia en 2025 no es el nombre de un proyecto en particular, sino que toda la industria llegó a un consenso más claro sobre la optimización a nivel de arquitectura: la separación PD, el EP grande y la colaboración nube-borde nos indican que la inferencia ya evolucionó de «hacer que funcione en una sola tarjeta» a «reconstruir el sistema según las características del workload».
